% tex-fmt: off
\documentclass[a4paper,11pt]{article} % Don't change this
\usepackage{fullpage} % Don't change this

\usepackage{titling} % Don't change this
\pretitle{\vspace{-6em}\begin{center}\larger[2]\bf} % Don't change this
\postauthor{\end{center}} % Don't change this
\preauthor{\vspace{-2em}\begin{center}\smaller[1]} % Don't change this
\postauthor{\end{center}\vspace{-5em}} % Don't change this

\usepackage[shortlabels]{enumitem} % Don't change this
\setlist{itemsep=0em,topsep=0.3em} % Don't change this
% tex-fmt: on

%%% Some useful packages.
\usepackage{xcolor}
\usepackage{float}
\usepackage{amsmath,amssymb,mathtools}
\usepackage{graphicx}
\usepackage{subcaption}
\usepackage[hidelinks]{hyperref}
\usepackage{relsize}
\usepackage{cleveref}
\crefname{figure}{Fig.}{Figs.}

\usepackage[natbib=true, style=authoryear-icomp, maxnames=3,
giveninits=true, uniquename=init]{biblatex}
\citetrackerfalse
\bibliography{ref} % put bib entries in ref.bib

\usepackage[linesnumbered]{algorithm2e}
\usepackage{cleveref}
\crefname{equation}{Equation.}{Equations.}
\crefname{figure}{Figure}{Figures.}

\usepackage{lipsum}
\usepackage{wrapfig}
\usepackage{tikz}

%%% Add additional packages or macros below
\usepackage[draft,layout=inline]{fixme}
\usepackage{microtype}
\usepackage{memoize}
\usepackage{pgfplots}
\usepackage{pgfplotstable}
\usepgfplotslibrary{groupplots}
\pgfplotsset{compat=1.18, table/col sep=comma}
\usepackage{pythontex}
\setpythontexworkingdir{<outputdir>}
\usepackage{siunitx}
\sisetup{drop-zero-decimal, round-mode=figures}
\usepackage{booktabs}
\crefname{table}{Table}{Tables}
\usepackage[section]{placeins}

\begin{document}

\begin{pycode}
import duckdb

database = "../../../results/grid.duckdb"
pytex.add_dependencies(database)
con = duckdb.connect(database, read_only=True)
for name in ("fits", "broken_line"):
    attached = f"../../../results/{name}.duckdb"
    pytex.add_dependencies(attached)
    con.sql(f"ATTACH '{attached}' AS {name} (READ_ONLY)")
query = con.sql("""
    SELECT *,
        (SELECT count(DISTINCT seed) FROM metrics NATURAL JOIN exemplar) AS n_seeds
    FROM exemplar
""")
cell = dict(zip(query.columns, query.fetchone()))
\end{pycode}

\title{Scale or Shape? Decomposing NTK Movement During Grokking}
\author{Group G2 \\
  Zain Al-Saffi (s4800836) $\cdot$
  Jasper Chong (s4883420) $\cdot$
  Aaron D'Mello (s4890257) \\
Zac Kienzle (s4803349) $\cdot$ Samuel Rieger (s4904959)}

\date{}
\maketitle

\begin{abstract}
  \fxfatal{Abstract, written last from the finished sections.}
\end{abstract}

\section{Introduction}

Grokking is the delayed generalisation that
\citet{powerGrokkingGeneralizationOverfitting2022} found on small algorithmic
datasets: validation accuracy rises from chance long after the training data
have been fitted. Modular addition learned by a two-layer perceptron is a
minimal instance. Full-batch gradient descent on the squared error groks there
with no explicit regularisation \citep{gromovGrokkingModularArithmetic2023},
and \citet{kumarGrokkingTransitionLazy2024} use this setting, with no weight
decay, as a baseline.

One account holds that grokking is a transition from lazy training, in which
the network behaves as a kernel method with its initial tangent kernel, to rich
training, in which the kernel adapts to the task
\citep{kumarGrokkingTransitionLazy2024}. The evidence that the kernel changes
is a distance: \citet{mohamadiWhyYouGrok2024} observe that the empirical NTK
stays close to its initial value until the training data are fitted and moves
far from it afterwards. A distance does not say how the kernel changed.
\citet{shanTheoryNeuralTangent2022} find that the tangent kernel grows in
Frobenius norm during training, as
\citet{baratinImplicitRegularizationNeural2021} report, but that kernel
gradient descent with the initial kernel scaled up to
the amplitude of the final one still trains more slowly than the network, so
that changes to the structure of the kernel, and not its amplitude alone,
account for the faster training.

The distinction matters for accounts of what sets the grokking time. The result
cited for weight decay forcing the NTK to change
\citep{kumarGrokkingTransitionLazy2024} is that of
\citet{lewkowyczTrainingDynamicsDeep2021}, in which the kernel only rescales,
and recent accounts time grokking by the weight norm
\citep{khanhWeightNormSets2026}, the scale of the logits
\citep{khanhWhatDoesWeight2026} and the rate of weight decay
\citep{kimGrokkingWeightDecayClock2026}. None of them measures the shape of the
kernel.

This report separates the two with published measures of the tangent kernel
(\cref{sec:statistics}): its scale $S_t$, the Frobenius norm relative to
initialisation; its shape change $R_t$, the kernel distance of
\citet{fortDeepLearningKernel2020}, which compares kernels in a scale-invariant
manner; its relative variation $D_t$ of
\citet{geigerDisentanglingFeatureLazy2020}; and its alignment $A_t$ with the
labels. We measure these statistics and the kernel's spectrum along training on
the modular-addition baseline of
\citet{kumarGrokkingTransitionLazy2024}, and ask which of scale and shape moves
before and while the network groks, and whether the times at which it memorises
and groks follow either.
\fxfatal{State the findings and the contributions
once the grid has run.}

\section{Related Work}

\paragraph{Grokking.} \Citet{powerGrokkingGeneralizationOverfitting2022}
report that on small algorithmic datasets, long after a network has overfitted
its training data, validation accuracy can rise from chance towards perfect
generalisation, and that weight decay is particularly effective at improving
generalisation there. \Citet{nandaProgressMeasuresGrokking2023} reverse-engineer
the algorithm a small transformer learns for modular addition, which turns
addition into rotation about a circle through discrete Fourier transforms, and
use it to split training into memorisation, circuit formation and clean-up.
For the two-layer perceptron we study,
\citet{gromovGrokkingModularArithmetic2023}
gives closed-form weights that solve modular addition when the activation is
quadratic, and observes grokking under full-batch gradient descent with no
explicit regularisation. \Citet{liuOmnigrokGrokkingAlgorithmic2023} trace
grokking to a mismatch between the training and test losses as functions of
the weight norm.

\paragraph{Lazy and rich training.} In the infinite-width limit the NTK stays
constant during training, so the network trains as a kernel method
\citep{jacotNeuralTangentKernel2020}.
\Citet{chizatLazyTrainingDifferentiable2020}
show that this lazy regime is a matter of scale: multiplying a model by a
factor $\alpha$, with its initial output kept bounded, brings its training
closer to that of its linearisation as $\alpha$ grows, while the
mean-field factor $1/m$ of a two-layer network of
width $m$ keeps the dynamics non-degenerate as $m\to\infty$.
\Citet{kumarGrokkingTransitionLazy2024} propose that grokking is a transition
from lazy to rich, feature-learning dynamics, which arises when the top
eigenvectors of the initial NTK are misaligned with the labels, the dataset is
large enough to generalise from but not so large that the training loss tracks
the test loss, and training starts in the lazy regime.
\Citet{mohamadiWhyYouGrok2024} observe that the Frobenius distance of the
empirical NTK from its initial value is negligible until the training data are
fitted and large afterwards, and prove that with fewer than $Cp^2$ training
pairs, for a constant $C$ they do not give, no input-output
permutation-equivariant kernel method has expected loss below half that of the
zero predictor on modular addition. \Citet{lyuDichotomyEarlyLate2024} prove that
homogeneous networks trained with large initialisation and small weight decay
stay at a kernel predictor for a long time and then move sharply to a
min-norm or max-margin predictor. \Citet{atanasovNeuralNetworksKernel2021}
describe silent alignment, in which the kernel's eigenstructure evolves while
the kernel is small and only its overall scale grows afterwards, under small
initialisation and whitened data.

\paragraph{Kernel alignment and spectral bias.} On a fixed kernel, gradient
descent fits the components of the labels along the eigenvectors with large
eigenvalues first (\cref{sec:spectra}). Networks learn low frequencies first
\citep{rahamanSpectralBiasNeural2019}, which
\citet{basriConvergenceRateNeural2019} explain for a two-layer network on the
sphere through the eigenvalues of its NTK at each frequency, and
\citet{canatarSpectralBiasTaskModel2021} relate the generalisation of kernel
regression to the share of the target's power in the top eigenfunctions.
Outside the kernel regime, the top eigenfunctions of the NTK align with the
target during training \citep{kopitkovNeuralSpectrumAlignment2020}, the NTK
changes rapidly in a short initial transient and at constant velocity afterwards
\citep{fortDeepLearningKernel2020}, and the spectrum of the tangent kernel
becomes more anisotropic while its centred alignment with the class labels
rises \citep{baratinImplicitRegularizationNeural2021}.
\Citet{shanTheoryNeuralTangent2022} give a theory of this alignment for deep
linear and two-layer ReLU networks, in which, with several outputs, the kernel
of each output aligns with its own target.
\Citet{kumarGrokkingTransitionLazy2024} also compute an alignment of the
initial kernel with the labels, on the test set; their printed formula has no
centring, and they evaluate it at initialisation only.

\paragraph{What sets the grokking time.} \Citet{khanhWeightNormSets2026} clamp
the weight norm at a multiple $\rho$ of a critical norm and find that the time
to grok grows exponentially in $\rho$. \Citet{khanhWhatDoesWeight2026} show that
under cross-entropy the norm acts on this delay through the scale of the logits,
while under squared error its effect is about half as large and runs through a
different route. \Citet{kimGrokkingWeightDecayClock2026} derive, for linear
models trained with heavy-ball momentum $\beta$ and weight decay, a
grokking time
that scales as $(1-\beta)/(\eta\lambda)$, and observe this scaling on modular
addition. \Citet{lewkowyczTrainingDynamicsDeep2021} prove, assuming a
conjecture on the
correlation functions of wide networks, that for a $k$-homogeneous network in
the infinite-width limit trained by gradient flow, $L_2$ regularisation makes
the NTK evolve as $\frac{d}{dt}\Theta_t=-2(k-1)\lambda\Theta_t$, and that under
squared error its eigenvectors stay fixed while its eigenvalues decay;
\citet{kumarGrokkingTransitionLazy2024} cite this result, in an argument they
call speculative, for weight decay moving networks out of the lazy regime.
The evolution in that result is a rescaling of the kernel. Its proof starts
from the gradient-flow equation of the kernel at finite width,
\begin{equation}
  \frac{d\Theta_t(x,x')}{dt} = -2(k-1)\lambda\Theta_t(x,x')
  + T_t(x,x') + T_t(x',x),
  \label{eq:kernelflow}
\end{equation}
where $T_t(x,x')=-\sum_{x''}\partial_{\mu\nu}f_t(x)\,\partial_\mu f_t(x')\,
\partial_\nu f_t(x'')\,\ell'(x'')$ sums over the training inputs and
repeated parameter indices, and shows that $T_0$ has expectation of order
$n^{-1}$ in the width $n$, so that $T$ drops out in the limit
\citep{lewkowyczTrainingDynamicsDeep2021}. In \eqref{eq:kernelflow} the
weight decay appears explicitly only in the term that rescales the kernel;
$T_t$ depends on it only through the trajectory. For homogeneous
networks trained on the squared error with weight decay,
\citet{pracherSpectralTheoryGrokking2026} show that a residual persists
after the training data are fitted, write the
evolution of the kernel as a term driven by that residual plus a weight-decay
term that shrinks the kernel, and predict a grokking time inversely
proportional to the product of learning rate and weight decay, which they
observe on modular addition; \citet{xuGrokGrokkingProvable2026} prove
bounds on the grokking time of over-parameterised linear regression under
gradient descent with weight decay. Both mechanisms need weight decay, and none
of these accounts separates the scale of the kernel from its shape.

\section{Methods}

\subsection{Task, network and training}

The task is the modular addition of \citet{gromovGrokkingModularArithmetic2023}
at the baseline of \citet{kumarGrokkingTransitionLazy2024}. For a modulus $p$,
each pair $(a,b)$ of residues in $\{0,\dots,p-1\}$ is encoded as the
concatenation $x_{a,b}\in\mathbb{R}^{2p}$ of the one-hot codes of $a$ and $b$,
and its target $y_{a,b}\in\mathbb{R}^{p}$ is the one-hot code of
$(a+b) \bmod p$. The dataset holds all $p^2$ pairs. Each seed draws its own
split of the pairs into a training set $\mathcal{T}$ and a test set, as well as
its own initial weights.

The network is the two-layer perceptron without biases of
\citet{gromovGrokkingModularArithmetic2023},
\begin{equation}
  f(x;\theta) = c_N W^{(2)} \phi\bigl(D^{-1/2} W^{(1)} x\bigr),
  \label{eq:network}
\end{equation}
with input dimension $D=2p$, width $N$, weights
$W^{(1)}\in\mathbb{R}^{N\times D}$ and $W^{(2)}\in\mathbb{R}^{p\times N}$
whose entries are drawn from $\mathcal{N}(0,1)$, and an activation $\phi$
applied to each entry. \Citet{gromovGrokkingModularArithmetic2023} derives
his exact solution for the quadratic activation, finds his empirical results
only mildly sensitive to the choice of activation, and reports that the same
weights work for the rectifier at a larger width; we use the rectifier
$\phi(h)=\max\{h,0\}$. A network function is $k$-homogeneous if scaling every
parameter by $c>0$ scales it by $c^k$, and a fully connected network with $L$
layers and rectifier activations is $L$-homogeneous
\citep{lewkowyczTrainingDynamicsDeep2021}, so \eqref{eq:network} is
2-homogeneous; the centred predictor below, which subtracts the constant
$f(x;\theta_0)$, is not. The readout factor is
$c_N=N^{-1/2}$ in the NTK parameterisation of
\citet{jacotNeuralTangentKernel2020} and $c_N=N^{-1}$ in the mean-field
parameterisation of \citet{gromovGrokkingModularArithmetic2023}. Following
\citet{kumarGrokkingTransitionLazy2024}, the trained predictor is the centred
and rescaled function
\begin{equation}
  \tilde f(x;\theta) = \alpha\bigl[f(x;\theta)-f(x;\theta_0)\bigr],
\end{equation}
where $\theta_0$ holds the initial weights, and the learning rate is
$\eta=\eta_0/\alpha^2$, which keeps the initial rate of change of the function
the same for every $\alpha$. The loss is the squared error averaged over the
$p$ outputs and the training pairs,
$L(\theta)=\frac{1}{p\lvert\mathcal{T}\rvert}\sum_{(a,b)\in\mathcal{T}}
\lVert\tilde f(x_{a,b};\theta)-y_{a,b}\rVert^2$, and training is full-batch
gradient descent with weight decay,
\begin{equation}
  \theta_{k+1} = (1-\eta\lambda)\,\theta_k - \eta\nabla L(\theta_k).
\end{equation}
This is the decoupled weight decay of
\citet{kimGrokkingWeightDecayClock2026} with no momentum buffer, and for plain
gradient descent coupled $L_2$ regularisation and decoupled weight decay
coincide \citep{kimGrokkingWeightDecayClock2026}. The grid sets the product
$\eta\lambda$ directly.

\subsection{The empirical tangent kernel}

For a scalar output $g(x;\theta)$ the tangent features are the gradients
$\nabla_\theta g(x;\theta)$, and the tangent kernel is the inner product of the
features of two inputs
\citep{jacotNeuralTangentKernel2020,baratinImplicitRegularizationNeural2021}.
We evaluate it on all $n=p^2$ pairs, so that at step $t$ the kernel is the
matrix $K_t\in\mathbb{R}^{n\times n}$ with entries
\begin{equation}
  [K_t]_{ij} = \bigl\langle\nabla_\theta g(x_i;\theta_t),
  \nabla_\theta g(x_j;\theta_t)\bigr\rangle .
\end{equation}
Two scalar outputs are measured: the sum of the logits
of \eqref{eq:network} divided by $\sqrt{p}$, whose kernel is the pseudo-NTK of
\citet{mohamadiFastWellFoundedApproximation2022} and is reported in the main
text, and the first logit, which \citet{mohamadiWhyYouGrok2024} use in place of
the full kernel of a vector output. The kernel is evaluated at
initialisation, on a
logarithmically spaced grid of steps, and at every event step of
\cref{sec:timing}.

\subsection{Scale, shape and alignment}
\label{sec:statistics}

For a set $\mathcal{S}$ of pairs, $K_t[\mathcal{S}]$ is the submatrix of $K_t$
on those pairs and $Y[\mathcal{S}]$ the matrix whose rows are their targets;
$\mathcal{T}$ holds the training pairs of a seed and $\mathcal{V}$ its test
pairs. Following \citet{cortesAlgorithmsLearningKernels2024}, the Frobenius
product is $\langle A,B\rangle_F=\operatorname{tr}(A^\top B)$ with norm
$\lVert A\rVert_F=\sqrt{\langle A,A\rangle_F}$, and $K_c=CKC$ centres an
$r\times r$ kernel matrix with $C=I_r-\frac{1}{r}\mathbf{1}\mathbf{1}^\top$.
All four statistics are measured on the test pairs of a seed. The first three
take the kernel $K_t=K_t[\mathcal{V}]$ of the summed logits; the alignment takes
the tangent kernel $\Theta_t$ of all $c$ logits on the same $n$ pairs, an
$nc\times nc$ matrix, and the vector $y\in\{0,1\}^{nc}$ that stacks their
one-hot targets $Y[\mathcal{V}]$ \citep[supplement, Sampled
Versions]{baratinImplicitRegularizationNeural2021}. We record
\begin{equation}
  \begin{aligned}
    S_t &= \frac{\lVert K_t\rVert_F}{\lVert K_0\rVert_F}, &
    R_t &= 1-\frac{\operatorname{tr}\bigl(K_tK_0^\top\bigr)}
    {\sqrt{\operatorname{tr}\bigl(K_tK_t^\top\bigr)}\,
    \sqrt{\operatorname{tr}\bigl(K_0K_0^\top\bigr)}},\\
    D_t &= \frac{\lVert K_t-K_0\rVert_F}{\lVert K_0\rVert_F}, &
    A_t &= \mathrm{CKA}\bigl(\Theta_t,yy^\top\bigr),\quad
    \mathrm{CKA}(K,K') = \frac{\operatorname{tr}(K_c{K'}_c)}
    {\lVert K_c\rVert_F\,\lVert {K'}_c\rVert_F}.
  \end{aligned}
  \label{eq:statistics}
\end{equation}
The scale $S_t$ is the Frobenius norm of the kernel relative to its value at
initialisation. \Citet{shanTheoryNeuralTangent2022} find that this norm grows
during training, and in the silent alignment of
\citet{atanasovNeuralNetworksKernel2021} the kernel first aligns while small
and then evolves in scale only, $K(x,x',t)=g(t)K_\infty(x,x')$. The shape
change $R_t$ is the kernel distance of \citet{fortDeepLearningKernel2020}
between the kernels at step $t$ and at initialisation, which compares two
kernel matrices in a scale-invariant manner. The variation $D_t$ is the total
relative variation of the kernel that \citet{geigerDisentanglingFeatureLazy2020}
measure on the test set to quantify feature training. The alignment $A_t$ is
the centred kernel alignment (CKA) of the tangent kernel with the class label
kernel $yy^\top$ on a held-out set, as
\citet{baratinImplicitRegularizationNeural2021}
measure it after \citet{cortesAlgorithmsLearningKernels2024}; its normalisation
makes it invariant under isotropic rescaling
\citep{baratinImplicitRegularizationNeural2021}. It takes all the logits
rather than their sum because the sum's kernel approximates the full one less
well as training proceeds \citep{mohamadiFastWellFoundedApproximation2022}.

Because the first three statistics are taken on one matrix, the variation is a
function of the scale and the shape,
\begin{equation}
  D_t^2 = S_t^2 - 2S_t(1-R_t) + 1 ,
  \label{eq:identity}
\end{equation}
and $R_t=0$ exactly when $K_t=cK_0$ for some $c>0$
(\cref{sec:identity}). A kernel that only rescales therefore has
$D_t=\lvert S_t-1\rvert$, which is as large as its change of scale, so a large
$D_t$ does not by itself show that the kernel changed shape.

\subsection{Spectral measures}
\label{sec:spectra}

In the linear dynamics of \citet{aroraFinegrainedAnalysisOptimization2019}, the
predictions $u_k$ on the training pairs follow
$u_{k+1}=u_k-\eta G(u_k-y)$ with a fixed Gram matrix $G$, from $u_0=0$. If $G$
has eigenvalues $\omega^G_i$ with orthonormal eigenvectors $g_i$, the residual
after $k$ steps is
$\lVert y-u_k\rVert=\sqrt{\sum_{i}(1-\eta\omega^G_i)^{2k}(g_i^\top y)^2}$, so
the part of the labels along an eigenvector with a large eigenvalue is fitted
first; \citet{jacotNeuralTangentKernel2020} give the continuous-time form, in
which the component along each kernel principal component decays exponentially
at a rate set by its eigenvalue. Eigenvalues are written $\omega$ to keep them
apart from the weight decay $\lambda$. The measures use the kernel on all
$n=p^2$ pairs, which covers every input of the task. The centred kernel
$CK_{t}C$, with $C=I_n-\frac{1}{n}\mathbf{1}\mathbf{1}^\top$, is symmetric and
positive semi-definite
\citep{cortesAlgorithmsLearningKernels2024}, so it has orthonormal eigenvectors
$v_1,\dots,v_n$ with eigenvalues $\omega_1\ge\dots\ge\omega_n\ge0$
\citep{axlerLinearAlgebraDone2024}. We record the target power along each
eigenvector, $P_{t,i}=\sum_{c=1}^{p}(v_i^\top Cy_c)^2$, where $y_c$ is the
$c$th column of $Y$, and its cumulative share over the leading eigenvectors,
which \citet{canatarSpectralBiasTaskModel2021} call the cumulative power
distribution and use to explain why a kernel generalises on one task and not on
another. The centring follows \citet{baratinImplicitRegularizationNeural2021},
who report results for centred tangent features;
\citet{jacotNeuralTangentKernel2020} note that their kernel principal
component analysis is not centred, so its first component is almost the
constant function. From the singular values
$\sigma_1\ge\dots\ge\sigma_n\ge0$ of $CK_{t}C$ we compute the effective rank of
\citet{royEffectiveRankMeasure2007}, which
\citet{baratinImplicitRegularizationNeural2021} apply to the tangent kernel, and
from the eigenvalues the trace ratios of
\citet{baratinImplicitRegularizationNeural2021},
\begin{equation}
  \operatorname{erank}_t = \exp\Bigl(-\sum_{j}\mu_j\log\mu_j\Bigr),\qquad
  \mu_j=\frac{\sigma_j}{\sum_{i}\sigma_i},\qquad
  T_{t,k}=\frac{\sum_{j\le k}\omega_j}{\sum_{j}\omega_j},
  \label{eq:erank}
\end{equation}
which measure how evenly the kernel spreads over its eigenvectors; $T_{t,k}$ is
the share of the trace carried by the $k$ largest eigenvalues, which is how
\citet{baratinImplicitRegularizationNeural2021} describe their trace ratios:
the relative importance of the top $k$ eigenvalues.
\Citet{royEffectiveRankMeasure2007} take $0\log0=0$, so
zero singular values, such as the one along the constant vector, which $C$ maps
to zero, do not enter the effective rank.

\subsection{Event times}
\label{sec:timing}

Following \citet{khanhWeightNormSets2026}, the memorisation step of a seed is
the first step at which its training accuracy reaches
\pys{\num{!{cell["memorised"]}}}, and its grokking step at level $\ell$ is the
first step at which its test accuracy reaches $\ell$, for each
$\ell\in\{\pys{\numlist{!{";".join(map(str, cell["grokked"]))}}}\}$, the
test-accuracy thresholds at which \citet{khanhWeightNormSets2026} check that
their critical weight norm does not depend on the threshold. A seed
that has not reached a level within the
step budget is kept, and its event time is recorded as censored at the budget.

Within a cell, the seeds give the Kaplan--Meier estimate of the probability
that an event has not occurred by step $t$,
\begin{equation}
  \hat F(t) = \prod_{j\mid t_j\le t}\frac{n_j-d_j}{n_j},
\end{equation}
where $t_1<t_2<\cdots$ are the distinct steps at which the event occurs, $d_j$
seeds reach it at $t_j$ and $n_j$ are still at risk there. Greenwood's formula
estimates its variance as
$\hat F^2(t)\sum_{j\mid t_j\le t}d_j/[n_j(n_j-d_j)]$, and with
$\hat s^2(t)=\sum_{j\mid t_j\le t}d_j/[n_j(n_j-d_j)]\big/[\log\hat F(t)]^2$
the approximate \qty{95}{\percent} interval $\hat F(t)^{\exp[\pm1.96\hat s(t)]}$
stays within $[0,1]$ \citep{kalbfleischStatisticalAnalysisFailure2002}. The
median event time of a cell is the step at which $\hat F$ falls to one half. A
cell whose $\hat F$ stays above one half up to the budget has no median within
it, and at a grokking level it does not grok.

Each run is also labelled by its phase at the last step, following
\citet{pracherSpectralTheoryGrokking2026}: grokking when its training and test
accuracies both end above a threshold, memorisation when only its training
accuracy does, forgetting when its training accuracy rose above the threshold
earlier and ends below it, and no fitting otherwise. The phases describe the
cells and do not select runs. A run keeps its event times whatever its phase,
since in the analysis of one failure mode a unit that fails by another mode, or
is removed from test without failing, contributes its running time as a
censoring time \citep{nelsonAcceleratedTestingStatistical1990}.

Three censored regressions, fitted with lifelines
\citep{davidson-pilonLifelinesSurvivalAnalysis2019}, relate the event times to
the settings of the cells. The log-normal accelerated failure time model takes
the survival function of an event time $T$ with covariates $x$ to be
\begin{equation}
  S(t\mid x) = 1-\Phi\biggl(\frac{\log t-\mu(x)}{\sigma}\biggr),\qquad
  \mu(x) = a_0 + \sum_{i}a_i x_i ,
  \label{eq:aft}
\end{equation}
where $\Phi$ is the standard normal distribution function, so a unit increase
in $x_i$ multiplies the median event time by $e^{a_i}$. Its parameters are
maximum likelihood estimates. With the step budget $\tau$ as the common
censoring time, seed $i$ contributes the log likelihood
\begin{equation}
  \mathcal{L}_i = I_i\Bigl[-\log\sigma-\tfrac{1}{2}\log(2\pi)
  -\tfrac{1}{2}z_i^2\Bigr] + (1-I_i)\log\bigl[1-\Phi(\zeta_i)\bigr]
\end{equation}
in terms of its log time, where $z_i=(\log t_i-\mu(x_i))/\sigma$,
$\zeta_i=(\log\tau-\mu(x_i))/\sigma$, and $I_i=1$ if the event is observed and
$I_i=0$ otherwise \citep{nelsonAcceleratedTestingStatistical1990}. This is the
likelihood of \citet{kalbfleischStatisticalAnalysisFailure2002} for a
log-linear model under independent censoring, with normal errors. The additive
model of
\citet{aalenLinearRegressionModel1989} takes the intensity of the event for
run $i$ at step $t$ to be $\alpha_0(t)+\sum_j\alpha_j(t)Z_j^i$ while the run
has not reached the level and zero afterwards, and estimates the cumulative
regression functions $A_j(t)=\int_0^t\alpha_j(s)\,ds$ by
\begin{equation}
  A^*(t) = \sum_{T_k\le t} \bigl[Y(T_k)^\top Y(T_k)\bigr]^{-1}Y(T_k)^\top I_k ,
  \label{eq:aalen}
\end{equation}
where the event times $T_1<T_2<\dots$ are ordered, row $i$ of $Y(t)$ is
$(1,Z_1^i,\dots,Z_r^i)$ while run $i$ is at risk and zero otherwise, and
$I_k$ is the indicator of the run whose event occurs at $T_k$. Unlike
\eqref{eq:aft} with logarithmic covariates, \eqref{eq:aalen} admits a covariate
at zero, such as $\eta\lambda=0$. The runs of one seed share its split and
initial weights, so the intervals resample seeds, each with all its runs and
their censoring, which keeps every group intact as the bootstrap for grouped
and censored data does \citep{davisonBootstrapMethodsTheir1997}, and take the
bias-corrected and accelerated interval
\citep{efronIntroductionBootstrap1994}.

Each event and level is fitted separately. The covariates of \eqref{eq:aft}
and \eqref{eq:aalen} are the settings of a cell, fixed in time. As
$\log(\eta\lambda)$ is undefined at $\eta\lambda=0$, \eqref{eq:aft} is fitted
twice: to the cells without weight decay, with indicators of the levels of
$\alpha$, and to the cells with $\eta\lambda>0$, all at the $\alpha$ of the
pilot cell, with the single covariate $\log(\eta\lambda)$, while
\eqref{eq:aalen} takes the pilot cell and the cells of the second stage
together, $\eta\lambda=0$ included, with $\eta\lambda$ itself. Every seed has
one run in each of these cells, the balanced grouping for which resampling
whole groups is justified \citep[Sec.~3.8]{davisonBootstrapMethodsTheir1997},
while the scan trains a single seed. The estimator is defined
only while $Y(t)$ has full rank, and the step at which the rank falls is a
final censoring time \citep[Sec.~4.1]{aalenLinearRegressionModel1989}; with
indicators of $\alpha$ it would stop when the first group of runs at one
$\alpha$ has left the risk set.
\citet{aalenLinearRegressionModel1989} states his results for covariates fixed
in time, and lifelines fits \eqref{eq:aalen} only for such covariates. The
statistics of \cref{sec:statistics} change along training, so they enter a
third model, the relative risk model with time-dependent covariates of
\citet{kalbfleischStatisticalAnalysisFailure2002},
\begin{equation}
  h\bigl(t\mid Z(t)\bigr) = h_0(t)\exp\bigl[Z(t)^\top\gamma\bigr],
\end{equation}
whose coefficients maximise the partial likelihood
\begin{equation}
  L(\gamma) = \prod_{j=1}^{k}
  \frac{\exp\bigl[Z_j(t_j)^\top\gamma\bigr]}
  {\sum_{l=1}^{n}Y_l(t_j)\exp\bigl[Z_l(t_j)^\top\gamma\bigr]} ,
\end{equation}
where $t_1<\cdots<t_k$ are the event times, $Z_j$ belongs to the run whose
event occurs at $t_j$, and $Y_l(t)$ is one while run $l$ is at risk.
\Citet{kalbfleischStatisticalAnalysisFailure2002} write it for untied event
times; runs can reach an event at the same step, and the lifelines fitter
handles such ties by Efron's method. Here
$Z(t)=(S_t,R_t,A_t)$ is held from each checkpoint to the next, so its paths
are left-continuous as the model requires. These covariates are internal:
the run that reaches the event generates them. The model therefore describes
the rate of the event given the current kernel and has no survivor function
of its own. When a setting acts on the event through these covariates, an
analysis conditional on them shows little effect of that setting
\citep{kalbfleischStatisticalAnalysisFailure2002}.

\subsection{Choice of levels}

Where weight decay sets the grokking time, \eqref{eq:aft} with the single
covariate $x=\log(\eta\lambda)$ makes $\log T$ normal with mean $a_0+a_1x$ and
variance $\sigma^2$. Without censoring this is the simple linear regression
model, whose designs \citet{bergerIntroductionOptimalDesigns2009} construct as
follows. A design $\xi$ places a proportion $w_j$ of the runs at level $x_j$,
$j=1,\dots,m$, with $\sum_j w_j=1$, on levels rescaled to $[-1,1]$, and has the
information matrix
\begin{equation}
  M(\xi) = \sum_{j=1}^{m} w_j f(x_j)f(x_j)^\top ,\qquad f(x) = (1,x)^\top .
\end{equation}
With $N$ runs, $n_j=Nw_j$ of them at $x_j$, the least squares estimates have
covariance $\sigma^2(X^\top X)^{-1}$, where $X^\top X=N M(\xi)$, and the slope
has variance $\operatorname{var}(\hat a_1)=\sigma^2/\sum_i(x_i-\bar x)^2$. The
design that minimises this variance is unique and places half the runs at each
end of the interval, and for the straight line it is also D-, A-, E- and
G-optimal. When the line is in doubt, a quadratic term measures its curvature.
The D-optimal design for the quadratic places a third of the runs at each of
$-1$, $0$ and $1$, and the design for its curvature coefficient alone places
half of them at $0$ and a quarter at each end. For many models with $p$
parameters an optimal design has between $p$ and $p(p+1)/2$ distinct levels,
and a design $\xi$ is compared with the D-optimal design $\xi^*$ of the model
that holds by its relative efficiency
\begin{equation}
  \mathrm{RE}_{\mathrm{D}} = \biggl\{
    \frac{\det\operatorname{Cov}(\hat a^*)}{\det\operatorname{Cov}(\hat a)}
  \biggr\}^{1/p} ,
\end{equation}
where $\hat a^*$ and $\hat a$ are the estimates under $\xi^*$ and $\xi$.

Censoring at the budget changes the information a run carries. For
$\mu(x)=a_0+a_1x$ and censoring at $\log\tau$, a run at level $x$ contributes
to the parameters $(a_0,a_1,\sigma)$ the expected information
\begin{equation}
  F_x = \frac{1}{\sigma^2}
  \begin{pmatrix}
    A(\zeta)  & xA(\zeta)   & B(\zeta)  \\
    xA(\zeta) & x^2A(\zeta) & xB(\zeta) \\
    B(\zeta)  & xB(\zeta)   & C(\zeta)
  \end{pmatrix},
\end{equation}
where $\zeta=(\log\tau-\mu(x))/\sigma$ and, with $\phi$ and $\Phi$ evaluated
at $\zeta$,
\begin{equation}
  \begin{aligned}
    A(\zeta) &= \Phi-\Bigl[\zeta-\frac{\phi}{1-\Phi}\Bigr]\phi, \qquad
    B(\zeta) = -\phi\Bigl[1+\zeta\Bigl(\zeta-\frac{\phi}{1-\Phi}\Bigr)\Bigr],\\
    C(\zeta) &= 2\Phi-\zeta\phi\Bigl[1+\zeta^2-\frac{\zeta\phi}{1-\Phi}\Bigr].
  \end{aligned}
\end{equation}
The information of a plan is the sum of $F_x$ over its runs, the covariance
of the maximum likelihood estimates is its inverse, and a function $h$ of the
parameters has variance $HF^{-1}H^\top$, where $H$ is the gradient of $h$
\citep{nelsonAcceleratedTestingStatistical1990}. Because $\zeta$ depends on the
parameters, an optimum plan for censored data is found by numerical search
for guessed values of them and is only locally optimum, and good compromise
plans run more units at low stress, here small $\eta\lambda$, than at high
\citep{nelsonAcceleratedTestingStatistical1990}.

Across the level at which the grokking time returns to its value without
weight decay, the mean log time is a broken line,
$\mu(x)=a_0+a_1x+b(x-\psi)_+$ with $(x-\psi)_+=(x-\psi)I(x>\psi)$. The method
of \citet{muggeoEstimatingRegressionModels2003} applies to any regression model
with a linear predictor, survival models included, so it runs inside the
censored likelihood of \eqref{eq:aft}. Step $s$ of the fit sets
$U^{(s)}=(x-\psi^{(s)})_+$ and $V^{(s)}=-I(x>\psi^{(s)})$, fits the linear
predictor $a_1x + b U^{(s)} + \kappa V^{(s)}$, and updates
$\psi^{(s+1)}=\psi^{(s)}+\hat\kappa/\hat b$ until $\hat\kappa$ vanishes, when
every estimate, $\hat\psi$ included, is maximum likelihood. Fixing the left
slope $a_1$ at zero gives a threshold model, in which the grokking time does
not change below the break, and whether the left slope is zero is a Wald test
of $a_1$ in the unconstrained fit \citep{muggeoSegmentedPackageFit2008}.
Whether a break exists at all is not a Wald test of $b$, since $\psi$
vanishes under $b=0$; the segmented package bounds the $p$-value of that test
by the method of Davies \citep{muggeoSegmentedPackageFit2008}. The break has
the Wald interval $\hat\psi\pm1.96\operatorname{SE}(\hat\psi)$, whose standard
error comes from the delta method for the ratio $\hat\kappa/\hat b$ and, at
convergence, where $\hat\kappa=0$, reduces to
$\operatorname{SE}(\hat\kappa)/\lvert\hat b\rvert$
\citep{muggeoEstimatingRegressionModels2003,muggeoSegmentedPackageFit2008}.
When the break lies away from the middle of the range of $x$, its estimate is
biased with an asymmetric distribution and the Wald interval can perform
poorly, in part because $\operatorname{SE}(\hat\psi)$ may understate its
variability \citep{muggeoIntervalEstimationBreakpoint2017}. The break
depends on constants that
the grid itself estimates, so a design for it is also only locally optimal; a
sequential design takes the guesses from a first stage, which here is a scan
of weight decays, and updates the design with their estimates
\citep{bergerIntroductionOptimalDesigns2009}.

\citet{parkExperimentalDesignsFitting1978a} constructs D-optimal designs for
segmented polynomials with a known break. For $k$ distinct levels,
$\lvert X^\top X\rvert$ is largest when every level has the same number of
runs, and under continuity at the break the design is D-optimal on each segment
with a level at the break. A line below the break and a quadratic above it
therefore have the design $x\in\{-1,\psi,(1+\psi)/2,1\}$ on the rescaled axis.
The quadratic is one degree above the line expected above the break, as
\citet{parkExperimentalDesignsFitting1978a} advises when lack of fit is to be
tested, with equal numbers of runs at each level. The design is locally optimal
at the estimated break. The first stage is the column of the grid of
\citet{pracherSpectralTheoryGrokking2026} at the baseline's learning rate, one
run at each weight decay, and its last grokking and first other point bound the
range above, as their boundary is placed. The broken line fitted to its runs
below that bound estimates the break, and the second stage places the design
at the estimate on the range from the column's lowest point to the bound.

\section{Experiments}

\subsection{The pilot cell}

The pilot cell is the baseline of \citet{kumarGrokkingTransitionLazy2024}:
modulus $p=\pys{\num{!{cell["p"]}}}$, a fraction
\pys{\num{!{cell["train_fraction"]}}} of the pairs for training, width
$N=\pys{\num{!{cell["width"]}}}$ in the
\pys{!{ {"ntk": "NTK", "mean_field": "mean-field"}[cell["parameterisation"]] }}
parameterisation, $\alpha=\pys{\num{!{cell["alpha"]}}}$,
$\eta_0=\pys{\num{!{cell["eta_0"]}}}$,
$\eta\lambda=\pys{\num{!{cell["eta_lambda"]}}}$ and a budget of
\pys{\num{!{cell["steps"]}}} steps, for \pys{!{cell["n_seeds"]}} seeds.
\Cref{tab:events} gives the steps at which the seeds memorise and grok,
\cref{fig:dynamics} the losses and the statistics \eqref{eq:statistics} along
training, and \cref{tab:kernel} those statistics and the spectral measures at
initialisation and at the last step.
\fxfatal{Describe the pilot against the questions of the introduction once the
grid gives the comparison cells.}

\begin{table}[t]
  \centering
  \caption{Steps at which the seeds of the pilot cell first reach each event:
    training accuracy at the memorisation level, or test accuracy at a grokking
    level. Median, minimum and maximum over seeds, and the number of seeds that
  reach the event within the budget.}
  \label{tab:events}
\begin{pycode}
events = con.sql("""
    SELECT event AS Event, level AS Level, median(time) AS Median,
        min(time) AS Minimum, max(time) AS Maximum,
        count(*) FILTER (WHERE observed) AS Reached
    FROM events NATURAL JOIN exemplar
    GROUP BY event, level
    ORDER BY event DESC, level
""").df()
print(events.style.hide(axis="index").to_latex(hrules=True, siunitx=True))
\end{pycode}
\end{table}

\begin{table}[t]
  \centering
  \caption{Statistics \eqref{eq:statistics}, and spectral measures
    \eqref{eq:erank} of the sum kernel, of the pilot cell at
    initialisation and at the last step, and the share of the centred target's
    power on the $2(p-1)$ leading eigenvectors. Minimum, median and maximum over
  seeds.}
  \label{tab:kernel}
\begin{pycode}
kernel = con.sql("""
    SELECT statistic AS Statistic, step AS Step, min(value) AS Minimum,
        median(value) AS Median, max(value) AS Maximum
    FROM (
        UNPIVOT (
            SELECT seed, step, S_sum, R_sum, A_full, D_sum
            FROM metrics NATURAL JOIN exemplar
            WHERE step IN (0, steps)
        ) ON S_sum, R_sum, A_full, D_sum INTO NAME statistic VALUE value
        UNION ALL BY NAME
        UNPIVOT (
            SELECT seed, step, erank, trace_ratio, target_share
            FROM spectral_summary NATURAL JOIN exemplar
            WHERE kernel = 'sum' AND step IN (0, steps)
        ) ON erank, trace_ratio, target_share INTO NAME statistic VALUE value
    )
    GROUP BY statistic, step
    ORDER BY statistic, step
""").df()
labels = {
    "S_sum": "$S_t$",
    "R_sum": "$R_t$",
    "A_full": "$A_t$",
    "D_sum": "$D_t$",
    "erank": r"$\operatorname{erank}_t$",
    "trace_ratio": "$T_{t,2(p-1)}$",
    "target_share": r"target power on $v_1,\dots,v_{2(p-1)}$",
}
print(kernel.style.hide(axis="index").format({"Statistic": labels.get}).to_latex(
    hrules=True, siunitx=True,
))
\end{pycode}
\end{table}

\subsection{The grid}

\fxfatal{The grid: its arms and levels, each level taken from a transcribed
source, and the timing fits.}

\begin{figure}[t]
  \centering
\begin{pycode}
series = "'^(train_loss|test_loss|S_sum|R_sum|A_full)$'"
con.sql(rf"""
    SELECT step,
        median(COLUMNS({series})) AS '\0',
        min(COLUMNS({series})) AS '\0_lo',
        max(COLUMNS({series})) AS '\0_hi'
    FROM metrics NATURAL JOIN exemplar
    WHERE step > 0
    GROUP BY step
    ORDER BY step
""").to_csv("dynamics.csv")
pytex.add_created("dynamics.csv")
print(r"\pgfplotstableread{pythontex-files-\jobname/dynamics.csv}\dynamics
\mmznext{
  context={%
    \csname pdf@filemdfivesum\endcsname{pythontex-files-\jobname/dynamics.csv}%
  },
}
\begin{tikzpicture}
  \begin{groupplot}[
      group style={
        group size=2 by 2,
        xlabels at=edge bottom,
        xticklabels at=edge bottom,
      },
      width=\linewidth/2,
      xmode=log,
      xlabel={step $t$},
      error bars/y dir=both,
      error bars/y explicit,
    ]
    \nextgroupplot[
      ymode=log,
      title={MSE},
      legend entries={train, test},
      legend pos=south west,
    ]
    \addplot table [
      x=step, y=train_loss,
      y error plus expr=\thisrow{train_loss_hi}-\thisrow{train_loss},
      y error minus expr=\thisrow{train_loss}-\thisrow{train_loss_lo},
    ] {\dynamics};
    \addplot table [
      x=step, y=test_loss,
      y error plus expr=\thisrow{test_loss_hi}-\thisrow{test_loss},
      y error minus expr=\thisrow{test_loss}-\thisrow{test_loss_lo},
    ] {\dynamics};
    \nextgroupplot[title={$S_t$}]
    \addplot table [
      x=step, y=S_sum,
      y error plus expr=\thisrow{S_sum_hi}-\thisrow{S_sum},
      y error minus expr=\thisrow{S_sum}-\thisrow{S_sum_lo},
    ] {\dynamics};
    \nextgroupplot[title={$R_t$}]
    \addplot table [
      x=step, y=R_sum,
      y error plus expr=\thisrow{R_sum_hi}-\thisrow{R_sum},
      y error minus expr=\thisrow{R_sum}-\thisrow{R_sum_lo},
    ] {\dynamics};
    \nextgroupplot[title={$A_t$}]
    \addplot table [
      x=step, y=A_full,
      y error plus expr=\thisrow{A_full_hi}-\thisrow{A_full},
      y error minus expr=\thisrow{A_full}-\thisrow{A_full_lo},
    ] {\dynamics};
  \end{groupplot}
\end{tikzpicture}
")
\end{pycode}
  \caption{Train and test MSE, the statistics $S_t$ and $R_t$ of the
    sum-over-logits kernel and the alignment $A_t$ of the tangent kernel of all
    logits \eqref{eq:statistics} against step $t$, at the baseline cell of
    \citet{kumarGrokkingTransitionLazy2024}. Points are medians over
  \py{cell["n_seeds"]} seeds and error bars span their range.}
  \label{fig:dynamics}
\end{figure}

\begin{figure}[t]
  \centering
\begin{pycode}
con.sql("""
    SELECT alpha, step, median(test_acc) AS test_acc
    FROM metrics NATURAL JOIN cells
    WHERE stage = '1' AND step > 0
    GROUP BY ALL
    ORDER BY alpha, step
""").to_csv("arms.csv")
pytex.add_created("arms.csv")
arms = [a for (a,) in con.sql(
    "SELECT alpha FROM cells WHERE stage = '1' ORDER BY alpha"
).fetchall()]
print(r"\def\armvalues{" + ",".join(map(str, arms)) + "}")
print(r"\pgfplotstableread{pythontex-files-\jobname/arms.csv}\arms
\mmznext{
  context={%
    \csname pdf@filemdfivesum\endcsname{pythontex-files-\jobname/arms.csv}%
  },
}
\begin{tikzpicture}
  \begin{axis}[
      xmode=log,
      xlabel={step $t$},
      ylabel={test accuracy},
      legend pos=outer north east,
    ]
    \expandafter\pgfplotsinvokeforeach\expandafter{\armvalues}{
      \addplot+[mark=none, restrict expr to domain={\thisrow{alpha}}{#1:#1}]
      table [x=step, y=test_acc] {\arms};
      \addlegendentry{$\alpha=#1$}
    }
  \end{axis}
\end{tikzpicture}
")
\end{pycode}
  \caption{Test accuracy against step $t$ for each output scale $\alpha$ of
    the first stage, the sweep of
    \citet[Fig.~2c]{kumarGrokkingTransitionLazy2024},
  as the median over \py{cell["n_seeds"]} seeds.}
  \label{fig:arms}
\end{figure}

\begin{table}[t]
  \centering
  \caption{The sequential design
    \citep[Sec.~5.4]{bergerIntroductionOptimalDesigns2009}. The scan's last
    grokking point, its first other point and their geometric midpoint, the
    upper end of the range
    \citep[App.~G.2.1]{pracherSpectralTheoryGrokking2026};
    the join $\hat\psi$ of the broken line fitted to the scan runs, with its
    interval \citep{muggeoSegmentedPackageFit2008}; the lower end of the scan;
    and the four levels of \citet[Sec.~4.1]{parkExperimentalDesignsFitting1978a}
  at $x$, $\ln\eta\lambda$ rescaled to $[-1,1]$.}
  \label{tab:design}
\begin{pycode}
design = con.sql(r"""
    UNPIVOT (
        SELECT b.last AS "last grokking scan point",
            b.first AS "first other scan point",
            b.upper AS "upper end",
            exp(j.psi_low) AS "join, lower limit",
            exp(j.psi) AS "join $\hat\psi$",
            exp(j.psi_high) AS "join, upper limit",
            d.lower AS "lower end"
        FROM boundary b, join_estimate j, design d
    ) ON COLUMNS(*) INTO NAME Quantity VALUE "$\eta\lambda$"
    UNION ALL BY NAME
    SELECT 'level ' || row_number() OVER (ORDER BY x) AS Quantity,
        eta_lambda AS "$\eta\lambda$", x AS "$x$"
    FROM design_levels
""").df()
print(design.style.hide(axis="index").format(na_rep="").to_latex(
    hrules=True, siunitx=True,
))
\end{pycode}
\end{table}

\begin{figure}[t]
  \centering
\begin{pycode}
con.sql("""
    SELECT eta_lambda, phase
    FROM phases NATURAL JOIN cells
    WHERE stage = 'scan'
    ORDER BY eta_lambda
""").to_csv("phases.csv")
pytex.add_created("phases.csv")
outcomes = [p for (p,) in con.sql(
    "SELECT DISTINCT phase FROM phases ORDER BY phase"
).fetchall()]
print(r"\def\phasenames{" + ",".join(outcomes) + "}")
print(r"\def\boundaryupper{" + str(con.sql("SELECT upper FROM boundary").fetchone()[0]) + "}")
print(r"\pgfplotstableread{pythontex-files-\jobname/phases.csv}\phases
\mmznext{
  context={%
    \csname pdf@filemdfivesum\endcsname{pythontex-files-\jobname/phases.csv}%
  },
}
\begin{tikzpicture}
  \begin{axis}[
      xmode=log,
      xlabel={$\eta\lambda$},
      symbolic y coords/.expand once=\phasenames,
      ytick=data,
      extra x ticks/.expand once=\boundaryupper,
      extra x tick labels={},
      extra x tick style={grid=major, major grid style={dashed}},
    ]
    \addplot+[only marks] table [x=eta_lambda, y=phase] {\phases};
  \end{axis}
\end{tikzpicture}
")
\end{pycode}
  \caption{Outcome of each scan run at the last step against $\eta\lambda$,
    the phases of \citet[Fig.~3a, App.~G.2]{pracherSpectralTheoryGrokking2026}
    along one column of their optimiser plane. The dashed line is the upper end
  of the range in \cref{tab:design}.}
  \label{fig:phases}
\end{figure}

\begin{figure}[t]
  \centering
\begin{pycode}
panels = {
    "memorised": ("memorised", cell["memorised"]),
    "grokked": ("grokked", cell["threshold"]),
}
for name, (event, level) in panels.items():
    con.sql("""
        SELECT ln(eta_lambda) AS x, ln(time) AS y, observed::INTEGER AS observed
        FROM design_runs
        WHERE event = $event AND level = $level
        ORDER BY x
    """, params={"event": event, "level": level}).to_csv(f"broken-{name}.csv")
    pytex.add_created(f"broken-{name}.csv")
    [fit] = con.sql("""
        SELECT a0, a1, b, psi, psi_low, psi_high, x_min, x_max, model
        FROM broken_line.broken_line
        WHERE event = $event AND level = $level
        QUALIFY row_number() OVER (ORDER BY aic) = 1
    """, params={"event": event, "level": level}).fetchall()
    for key, value in zip(
        ("A", "B", "U", "Psi", "Low", "High", "Min", "Max", "Model"), fit
    ):
        print(rf"\def\{name}{key}{{{value}}}")
    print(rf"\def\{name}Level{{{level}}}")
print(r"\mmznext{
  context={%
    \csname
    pdf@filemdfivesum\endcsname{pythontex-files-\jobname/broken-memorised.csv}%
    \csname
    pdf@filemdfivesum\endcsname{pythontex-files-\jobname/broken-grokked.csv}%
  },
}
\begin{tikzpicture}
  \begin{groupplot}[
      group style={group size=2 by 1, ylabels at=edge left},
      width=\linewidth/2,
      xlabel={$\ln\eta\lambda$},
      ylabel={$\ln T$},
      scatter/classes={0={mark=o}, 1={mark=*}},
    ]
    \nextgroupplot[title={memorised, level \memorisedLevel}]
    \addplot[scatter, only marks, scatter src=explicit symbolic]
    table [x=x, y=y, meta=observed]
    {pythontex-files-\jobname/broken-memorised.csv};
    \addplot[samples at={\memorisedMin,\memorisedPsi,\memorisedMax}, thick]
    {\memorisedA + \memorisedB*x + \memorisedU*max(x-\memorisedPsi, 0)};
    \draw[|-|] ({axis cs:\memorisedLow,0}|-{rel axis cs:0,0.05})
    -- ({axis cs:\memorisedHigh,0}|-{rel axis cs:0,0.05});
    \nextgroupplot[title={grokked, level \grokkedLevel}]
    \addplot[scatter, only marks, scatter src=explicit symbolic]
    table [x=x, y=y, meta=observed]
    {pythontex-files-\jobname/broken-grokked.csv};
    \addplot[samples at={\grokkedMin,\grokkedPsi,\grokkedMax}, thick]
    {\grokkedA + \grokkedB*x + \grokkedU*max(x-\grokkedPsi, 0)};
    \draw[|-|] ({axis cs:\grokkedLow,0}|-{rel axis cs:0,0.05})
    -- ({axis cs:\grokkedHigh,0}|-{rel axis cs:0,0.05});
  \end{groupplot}
\end{tikzpicture}
")
\end{pycode}
  \caption{Log event time against $\ln\eta\lambda$ for the runs of the design,
    with the broken line of \citet{muggeoEstimatingRegressionModels2003} chosen
    by AIC between a free and a zero left slope, and its join $\hat\psi$ with
    the delta-method interval drawn as a bar, as in
    \citet[Fig.~1]{muggeoSegmentedPackageFit2008}. Filled marks are observed
  events and open marks runs censored at the budget.}
  \label{fig:broken}
\end{figure}

\begin{figure}[t]
  \centering
\begin{pycode}
con.sql("""
    SELECT c.cell, c.stage, s.time, s.survival, s.lower, s.upper
    FROM fits.survival s JOIN cells c USING (cell)
    WHERE s.event = 'grokked'
        AND s.level = (SELECT threshold FROM exemplar)
        AND c.stage IN ('1', '2')
    ORDER BY c.cell, s.time
""").to_csv("survival.csv")
pytex.add_created("survival.csv")
arms = con.sql(
    "SELECT cell, alpha FROM cells WHERE stage = '1' ORDER BY alpha"
).fetchall()
levels = con.sql(
    "SELECT cell FROM cells WHERE stage = '2' ORDER BY eta_lambda"
).fetchall()
print(r"\def\armcells{" + ",".join(str(c) for c, _ in arms) + "}")
print(r"\def\armlegend{" + ",".join(rf"{{$\alpha={a}$}}" for _, a in arms) + "}")
print(r"\def\levelcells{" + ",".join(str(c) for (c,) in levels) + "}")
print(r"\def\levellegend{"
      + ",".join(f"{{level {k}}}" for k in range(1, len(levels) + 1)) + "}")
print(r"\pgfplotstableread{pythontex-files-\jobname/survival.csv}\survival
\mmznext{
  context={%
    \csname pdf@filemdfivesum\endcsname{pythontex-files-\jobname/survival.csv}%
  },
}
\begin{tikzpicture}
  \begin{groupplot}[
      group style={group size=2 by 1, ylabels at=edge left},
      width=\linewidth/2,
      xlabel={step $t$},
      ylabel={$\hat F(t)$},
      ymin=0, ymax=1,
    ]
    \nextgroupplot[title={first stage},
    legend entries/.expand once=\armlegend, legend to name=survivalarms]
    \expandafter\pgfplotsinvokeforeach\expandafter{\armcells}{
      \addplot+[const plot, mark=none, dashed, forget plot,
      restrict expr to domain={\thisrow{cell}}{#1:#1}]
      table [x=time, y=lower] {\survival};
      \addplot+[const plot, mark=none, dashed, forget plot,
      restrict expr to domain={\thisrow{cell}}{#1:#1}]
      table [x=time, y=upper] {\survival};
      \addplot+[const plot, mark=none,
      restrict expr to domain={\thisrow{cell}}{#1:#1}]
      table [x=time, y=survival] {\survival};
    }
    \nextgroupplot[title={second stage},
    legend entries/.expand once=\levellegend, legend to name=survivallevels]
    \expandafter\pgfplotsinvokeforeach\expandafter{\levelcells}{
      \addplot+[const plot, mark=none, dashed, forget plot,
      restrict expr to domain={\thisrow{cell}}{#1:#1}]
      table [x=time, y=lower] {\survival};
      \addplot+[const plot, mark=none, dashed, forget plot,
      restrict expr to domain={\thisrow{cell}}{#1:#1}]
      table [x=time, y=upper] {\survival};
      \addplot+[const plot, mark=none,
      restrict expr to domain={\thisrow{cell}}{#1:#1}]
      table [x=time, y=survival] {\survival};
    }
  \end{groupplot}
\end{tikzpicture}

\ref{survivalarms}\quad\ref{survivallevels}
")
\end{pycode}
  \caption{Kaplan--Meier estimates of the probability of not yet having grokked
    at level \py{cell["threshold"]}, by output scale in the first stage and by
    weight decay level in the second, with Greenwood's exponential interval
  dashed \citep[Sec.~1.4.1]{kalbfleischStatisticalAnalysisFailure2002}.}
  \label{fig:survival}
\end{figure}

\begin{figure}[t]
  \centering
\begin{pycode}
con.sql("""
    SELECT rank, (step = 0)::INTEGER AS initial,
        median(cumulative_share) AS share
    FROM target_power NATURAL JOIN exemplar
    WHERE kernel = 'sum' AND step IN (0, steps)
    GROUP BY ALL
    ORDER BY initial DESC, rank
""").to_csv("spectrum.csv")
pytex.add_created("spectrum.csv")
print(rf"\def\leading{{{2 * (cell['p'] - 1)}}}")
print(r"\pgfplotstableread{pythontex-files-\jobname/spectrum.csv}\spectrum
\mmznext{
  context={%
    \csname pdf@filemdfivesum\endcsname{pythontex-files-\jobname/spectrum.csv}%
  },
}
\begin{tikzpicture}
  \begin{axis}[
      xmode=log,
      xlabel={rank $\rho$},
      ylabel={$C(\rho)$},
      ymin=0, ymax=1,
      legend pos=outer north east,
    ]
    \addplot+[mark=none, restrict expr to domain={\thisrow{initial}}{1:1}]
    table [x=rank, y=share] {\spectrum};
    \addlegendentry{initialisation}
    \addplot+[mark=none, restrict expr to domain={\thisrow{initial}}{0:0}]
    table [x=rank, y=share] {\spectrum};
    \addlegendentry{last step}
    \draw[dotted] ({axis cs:\leading,0}|-{rel axis cs:0,0})
    -- ({axis cs:\leading,0}|-{rel axis cs:0,1});
  \end{axis}
\end{tikzpicture}
")
\end{pycode}
  \caption{Cumulative share $C(\rho)$ of the centred target's power on the
    leading $\rho$ eigenvectors of the sum kernel of the pilot cell, at
    initialisation and at the last step, the curve of
    \citet[Eq.~6, Fig.~1e]{canatarSpectralBiasTaskModel2021}. Medians over
  \py{cell["n_seeds"]} seeds; the dotted line marks $\rho=2(p-1)$.}
  \label{fig:spectrum}
\end{figure}

\begin{table}[t]
  \centering
  \caption{Coefficients of the log-normal accelerated failure time model
    \eqref{eq:aft} with their \qty{95}{\percent} intervals, for the output
    scale arm against $\alpha=\py{cell["alpha"]}$ and the weight decay arm on
  $\ln\eta\lambda$.}
  \label{tab:aft}
\begin{pycode}
aft = con.sql(r"""
    SELECT arm AS Arm, event AS Event, level AS Level,
        CASE WHEN covariate LIKE 'C(alpha%'
            THEN '$\alpha=' || regexp_extract(covariate, 'T\.([0-9.]+)', 1) || '$'
            ELSE '$\ln\eta\lambda$' END AS Covariate,
        coef AS Coefficient, "coef lower 95%" AS Lower, "coef upper 95%" AS Upper
    FROM fits.aft
    WHERE param = 'mu_' AND covariate != 'Intercept'
    ORDER BY arm, event DESC, level, covariate
""").df()
print(aft.style.hide(axis="index").to_latex(hrules=True, siunitx=True))
\end{pycode}
\end{table}

\begin{table}[t]
  \centering
  \caption{Coefficients of the relative risk model with time-dependent
    covariates \citep[Eq.~6.14]{kalbfleischStatisticalAnalysisFailure2002} on
    $S_t$, $R_t$ and $A_t$ \eqref{eq:statistics}, with their
    \qty{95}{\percent} intervals
  and $p$-values.}
  \label{tab:cox}
\begin{pycode}
cox = con.sql(r"""
    SELECT event AS Event, level AS Level,
        '$' || left(covariate, 1) || '_t$' AS Covariate,
        coef AS Coefficient, "coef lower 95%" AS Lower,
        "coef upper 95%" AS Upper, p AS "$p$"
    FROM fits.cox
    ORDER BY event DESC, level, covariate
""").df()
print(cox.style.hide(axis="index").to_latex(hrules=True, siunitx=True))
\end{pycode}
\end{table}

\section{Conclusion}

\fxfatal{Conclusion: what the measurements show, and the limitations.}

\section*{Contribution}

\fxfatal{Contribution of each group member, from the group.}

\section*{Acknowledgement}

\fxfatal{Acknowledgement, written by the group.}

\appendix
\section{Derivations}

\subsection{Scale, shape and variation}
\label{sec:identity}

Throughout, $K_t$ and $K_0$ are real $r\times r$ kernel matrices on the test
pairs of a seed with $K_0\neq0$ and $K_t\neq0$, so that \eqref{eq:statistics}
is defined.

\paragraph{Results used.} On the linear maps between two finite-dimensional
inner product spaces, $\langle S,T\rangle=\operatorname{tr}(T^*S)$ is an inner
product, which with respect to orthonormal bases is the standard inner product
of the matrix entries, and its norm is the Frobenius norm
\citep{axlerLinearAlgebraDone2024}. For real matrices the adjoint is the
transpose, so this is the Frobenius product
$\langle A,B\rangle_F=\operatorname{tr}(A^\top B)$ of
\citet{cortesAlgorithmsLearningKernels2024}. Square matrices satisfy
$\operatorname{tr}(AB)=\operatorname{tr}(BA)$
\citep{axlerLinearAlgebraDone2024}. In an inner product space,
\begin{equation*}
  \lVert u-v\rVert^2 = \lVert u\rVert^2 + \lVert v\rVert^2
  - \langle u,v\rangle - \langle v,u\rangle ,
\end{equation*}
over the real numbers $\langle u,v\rangle=\langle v,u\rangle$, and
$\langle u,v\rangle=\lVert u\rVert\,\lVert v\rVert$ holds if and only if one of
$u$ and $v$ is a nonnegative real multiple of the other
\citep{axlerLinearAlgebraDone2024}.

\paragraph{The identity.} By $\operatorname{tr}(AB)=\operatorname{tr}(BA)$,
$\operatorname{tr}(K_tK_0^\top)=\operatorname{tr}(K_0^\top K_t)
=\langle K_0,K_t\rangle_F$, and likewise
$\operatorname{tr}(K_tK_t^\top)=\lVert K_t\rVert_F^2$, so
\eqref{eq:statistics} gives
\begin{equation*}
  1-R_t = \frac{\langle K_t,K_0\rangle_F}{\lVert K_t\rVert_F\,\lVert
  K_0\rVert_F}.
\end{equation*}
The expansion above with the symmetry of the real inner product gives
\begin{equation*}
  \lVert K_t-K_0\rVert_F^2
  = \lVert K_t\rVert_F^2 + \lVert K_0\rVert_F^2
  - 2\langle K_t,K_0\rangle_F .
\end{equation*}
Dividing by $\lVert K_0\rVert_F^2$, the left side is $D_t^2$, the first term
is $S_t^2$, and
\begin{equation*}
  \frac{\langle K_t,K_0\rangle_F}{\lVert K_0\rVert_F^2}
  = \frac{\lVert K_t\rVert_F}{\lVert K_0\rVert_F}\cdot
  \frac{\langle K_t,K_0\rangle_F}{\lVert K_t\rVert_F\,\lVert K_0\rVert_F}
  = S_t(1-R_t),
\end{equation*}
which gives \eqref{eq:identity}.

\paragraph{When the shape is unchanged.} By the expression for $1-R_t$ above,
$R_t=0$ exactly when
$\langle K_t,K_0\rangle_F=\lVert K_t\rVert_F\,\lVert K_0\rVert_F$, that is,
by the last result above, exactly when one of $K_t$ and $K_0$ is a
nonnegative real multiple of the other. Neither is zero, so the multiple is
positive, and $R_t=0$ exactly when $K_t=cK_0$ for some $c>0$. Then
$S_t=c$, and \eqref{eq:identity} gives $D_t^2=(S_t-1)^2$.

\section{Model checks}

\begin{figure}[t]
  \centering
\begin{pycode}
from scipy.stats import norm

plot = con.sql("""
    SELECT alpha, time, observed,
        row_number() OVER (PARTITION BY alpha ORDER BY time) AS i,
        count(*) OVER (PARTITION BY alpha) AS n
    FROM events NATURAL JOIN cells
    WHERE stage = '1' AND event = 'grokked'
        AND level = (SELECT threshold FROM exemplar)
""").df()
plot = plot[plot.observed]
plot["position"] = norm.ppf((plot.i - 0.5) / plot.n)
plot.to_csv("probability.csv", index=False)
pytex.add_created("probability.csv")
print(r"\def\armvalues{" + ",".join(map(str, sorted(plot.alpha.unique()))) + "}")
print(r"\pgfplotstableread{pythontex-files-\jobname/probability.csv}\probability
\mmznext{
  context={%
    \csname
    pdf@filemdfivesum\endcsname{pythontex-files-\jobname/probability.csv}%
  },
}
\begin{tikzpicture}
  \begin{axis}[
      xmode=log,
      xlabel={grokking step $T$},
      ylabel={$\Phi^{-1}\bigl((i-0.5)/n\bigr)$},
      legend pos=outer north east,
    ]
    \expandafter\pgfplotsinvokeforeach\expandafter{\armvalues}{
      \addplot+[only marks, restrict expr to domain={\thisrow{alpha}}{#1:#1}]
      table [x=time, y=position] {\probability};
      \addlegendentry{$\alpha=#1$}
    }
  \end{axis}
\end{tikzpicture}
")
\end{pycode}
  \caption{Log-normal probability plot of the grokking times at level
    \py{cell["threshold"]} for each output scale of the first stage, each
    observed time at the normal quantile of its plotting position
    $(i-0.5)/n$ and runs censored at the budget left out
    \citep[Ch.~3, Sec.~4.2]{nelsonAcceleratedTestingStatistical1990}. Parallel
  straight lines support the log-normal model \eqref{eq:aft}.}
  \label{fig:probability}
\end{figure}

\begin{figure}[t]
  \centering
\begin{pycode}
aalen = con.sql(r"""
    SELECT dense_rank() OVER (ORDER BY covariate) AS panel,
        CASE WHEN covariate LIKE 'C(alpha%'
            THEN '$\alpha=' || regexp_extract(covariate, 'T\.([0-9.]+)', 1) || '$'
            WHEN covariate = 'Intercept' THEN 'intercept'
            ELSE '$\eta\lambda$' END AS title,
        time, estimate, low, high
    FROM fits.aalen
    WHERE event = 'grokked' AND level = (SELECT threshold FROM exemplar)
    ORDER BY panel, time
""")
aalen.select("panel, time, estimate, low, high").to_csv("aalen.csv")
pytex.add_created("aalen.csv")
titles = aalen.select("panel, title").distinct().order("panel").fetchall()
for panel, title in titles:
    print(rf"\expandafter\def\csname aalentitle{panel}\endcsname{{{title}}}")
print(r"\def\aalenpanels{" + ",".join(str(p) for p, _ in titles) + "}")
print(rf"\def\aalencount{{{len(titles)}}}")
print(rf"\def\aalengroup{{{len(titles)} by 1}}")
print(r"\pgfplotstableread{pythontex-files-\jobname/aalen.csv}\aalen
\mmznext{
  context={%
    \csname pdf@filemdfivesum\endcsname{pythontex-files-\jobname/aalen.csv}%
  },
}
\begin{tikzpicture}
  \begin{groupplot}[
      group style={group size/.expand once=\aalengroup, xlabels at=edge bottom},
      width=\linewidth/\aalencount,
      xlabel={step $t$},
    ]
    \expandafter\pgfplotsinvokeforeach\expandafter{\aalenpanels}{
      \nextgroupplot[title={\csname aalentitle#1\endcsname}]
      \addplot+[const plot, mark=none, dashed, forget plot,
      restrict expr to domain={\thisrow{panel}}{#1:#1}]
      table [x=time, y=low] {\aalen};
      \addplot+[const plot, mark=none, dashed, forget plot,
      restrict expr to domain={\thisrow{panel}}{#1:#1}]
      table [x=time, y=high] {\aalen};
      \addplot+[const plot, mark=none,
      restrict expr to domain={\thisrow{panel}}{#1:#1}]
      table [x=time, y=estimate] {\aalen};
    }
  \end{groupplot}
\end{tikzpicture}
")
\end{pycode}
  \caption{Cumulative regression functions of
    \citet{aalenLinearRegressionModel1989}
    for grokking at level \py{cell["threshold"]}, with the bias-corrected and
    accelerated bootstrap interval over seeds dashed, drawn as in
    \citet[Fig.~7]{aalenLinearRegressionModel1989}. A slope that changes along
  $t$ is an effect that changes with time.}
  \label{fig:aalen}
\end{figure}

\smaller
\printbibliography

\newpage
\fxfatal{The group keeps one of the two consent statements below.}

\begin{center}
  \emph{We give consent for this to be used as a teaching resource.}\\
  \emph{We DO NOT give consent for this to be used as a teaching resource.}
\end{center}

\end{document}
